\section{第3章 微分中值定理与导数的应用}
\subsection{中值定理}
\begin{exercise}{289}{基础·罗尔定理条件}
判断 $f(x)=|x|$ 与 $g(x)=x^2-1$ 在 $[-1,1]$ 上是否满足罗尔定理条件；对满足条件者求罗尔点。
\end{exercise}
\begin{exercise}{290}{基础·罗尔点计算}
验证 $f(x)=\ln\sin x$ 在 $[\frac{\pi}{6},\frac{5\pi}{6}]$ 上满足罗尔定理，并求 $f'(\xi)=0$ 的点。
\end{exercise}
\begin{exercise}{291}{常规·辅助函数与罗尔定理}
设 $f\in C[0,1]$，在 $(0,1)$ 可导，$f(0)=0$、$f(1)=\frac{\pi}{4}$。对 $F(x)=f(x)-\arctan x$ 使用罗尔定理，证明存在 $\xi\in(0,1)$ 使 $(1+\xi^2)f'(\xi)=1$。
\end{exercise}
\begin{exercise}{292}{常规·多项式与罗尔定理}
设 $a+\frac{b}{2}+\frac{c}{3}=0$。对 $F(x)=ax+\frac{bx^2}{2}+\frac{cx^3}{3}$ 使用罗尔定理，证明 $a+bx+cx^2$ 在 $(0,1)$ 内有零点。
\end{exercise}
\begin{exercise}{293}{常规·两次罗尔定理}
设 $f$ 在包含 $[0,1]$ 的开区间二阶可导，$f(0)=f(1)=0$、$f'(0)=0$。证明存在 $\xi\in(0,1)$ 使 $f''(\xi)=0$。提示：先对 $f$ 使用罗尔定理，再对 $f'$ 使用。
\end{exercise}
\begin{exercise}{294}{常规·乘积辅助函数}
设 $f\in C^2[0,1]$ 且 $f(1)=0$，令 $F(x)=x^2f(x)$。求 $F(0)$、$F'(0)$，并用两次罗尔定理证明 $F''$ 在 $(0,1)$ 内有零点。
\end{exercise}
\begin{exercise}{295}{常规·指数辅助函数}
设 $f\in C[0,1]$，在 $(0,1)$ 可导，$f(0)=1$、$f(1)=e^{-1}$。对 $F(x)=f(x)-e^{-x}$ 使用罗尔定理，证明存在 $\xi\in(0,1)$ 使 $f'(\xi)+e^{-\xi}=0$。
\end{exercise}
\begin{exercise}{296}{常规·乘积函数零点}
设 $a>0$，$f\in C[0,a]$，在 $(0,a)$ 可导且 $f(a)=0$。对 $F(x)=xf(x)$ 使用罗尔定理，证明存在 $\xi\in(0,a)$ 使 $f(\xi)+\xi f'(\xi)=0$。
\end{exercise}
\begin{exercise}{297}{常规·加权导数零点}
设 $f\in C[0,1]$，在 $(0,1)$ 可导且 $f(1)=0$。对 $F(x)=x^2f(x)$ 使用罗尔定理，证明存在 $\xi\in(0,1)$ 使 $2f(\xi)+\xi f'(\xi)=0$。
\end{exercise}
\begin{exercise}{298}{基础·对称导数等式}
设 $f$ 在 $(0,1)$ 可导且处处非零。直接取 $\xi=\frac{1}{2}$，验证 $\frac{f'(\xi)}{f(\xi)}=\frac{f'(1-\xi)}{f(1-\xi)}$。
\end{exercise}
\begin{exercise}{299}{常规·平方辅助函数}
设 $f\in C[0,1]$，在 $(0,1)$ 可导，$f(0)=1$、$f(1)=\sqrt e$。对 $F(x)=f(x)^2-e^x$ 使用罗尔定理，证明存在 $\xi\in(0,1)$ 使 $2f(\xi)f'(\xi)=e^\xi$。
\end{exercise}
\begin{exercise}{300}{基础·驻点存在性}
设 $f\in C[0,3]$，在 $(0,3)$ 可导，$f(0)=f(3)$。说明存在 $\xi\in(0,3)$ 使 $f'(\xi)=0$，并写出所用定理。
\end{exercise}
\begin{exercise}{301}{常规·介值与中值定理}
设 $f\in C[0,1]$，在 $(0,1)$ 可导，$f(0)=f(1)=0$、$f(\frac{1}{2})=1$。先对 $f(x)-x$ 使用零点定理，证明存在 $\eta\in(\frac{1}{2},1)$ 使 $f(\eta)=\eta$；再在 $[0,\eta]$ 上证明存在 $\xi$ 使 $f'(\xi)=1$。
\end{exercise}
\begin{exercise}{302}{常规·三零点与二阶导数}
设 $f\in C^2[0,1]$，$f(0)=f(\frac{1}{2})=f(1)=0$。使用两次罗尔定理证明存在 $\xi\in(0,1)$ 使 $f''(\xi)=0$。
\end{exercise}
\begin{exercise}{303}{基础·拉格朗日中值定理}
分别判断 $\ln(1+x^2)$、$|x|$ 在 $[-1,1]$ 上是否满足拉格朗日中值定理条件；对前者求中值点。
\end{exercise}
\begin{exercise}{304}{基础·中值点计算}
对 $f(x)=x^2$ 在 $[1,3]$ 上使用拉格朗日中值定理，求满足 $f(3)-f(1)=2f'(\xi)$ 的 $\xi$。
\end{exercise}
\begin{exercise}{305}{常规·三角函数不等式}
用拉格朗日中值定理证明 $\sin x<x$ 对 $0<x<\frac{\pi}{2}$ 成立。提示：对 $\sin t$ 在 $[0,x]$ 上使用定理。
\end{exercise}
\begin{exercise}{306}{常规·对数增量估计}
设 $a>b>0$。用 $f(x)=\ln x$ 的拉格朗日中值定理证明 $\frac{a-b}{a}<\ln(\frac{a}{b})<\frac{a-b}{b}$。
\end{exercise}
\begin{exercise}{307}{常规·对数不等式}
用拉格朗日中值定理证明 $\frac{x}{1+x}<\ln(1+x)<x$（$x>0$）。提示：对 $\ln(1+t)$ 在 $[0,x]$ 上使用定理。
\end{exercise}
\begin{exercise}{308}{常规·指数不等式}
用拉格朗日中值定理证明 $e^x>1+x$（$x>0$）。
\end{exercise}
\begin{exercise}{309}{常规·导数为零的函数}
在 $(-1,1)$ 上求 $F(x)=\arcsin x+\arccos x$ 的导数，利用中值定理证明 $F(x)=\frac{\pi}{2}$。
\end{exercise}
\begin{exercise}{310}{常规·反三角恒等式}
在 $x>0$ 上用求导证明 $\arctan x+\arctan(\frac{1}{x})=\frac{\pi}{2}$。提示：先证明左侧为常数，再取 $x=1$。
\end{exercise}
\begin{exercise}{311}{常规·乘积函数中值等式}
设 $f\in C[a,b]$，在 $(a,b)$ 可导，$a<b$。对 $F(x)=xf(x)$ 使用拉格朗日中值定理，证明存在 $\xi\in(a,b)$ 使 $bf(b)-af(a)=(b-a)[f(\xi)+\xi f'(\xi)]$。
\end{exercise}
\begin{exercise}{312}{常规·导数上界与增量}
设 $f\in C[a,b]$，在 $(a,b)$ 可导且 $|f'(x)|\le M$。用拉格朗日中值定理证明 $|f(b)-f(a)|\le M(b-a)$。
\end{exercise}
\begin{exercise}{313}{基础·柯西中值点}
对 $f(x)=\sin x$、$g(x)=1+\cos x$ 在 $[0,\frac{\pi}{2}]$ 上使用柯西中值定理，求满足 $[f(\frac{\pi}{2})-f(0)]g'(\xi)=[g(\frac{\pi}{2})-g(0)]f'(\xi)$ 的 $\xi$。
\end{exercise}
\begin{exercise}{314}{常规·柯西中值定理}
设 $0<a<b$，$f\in C[a,b]$，在 $(a,b)$ 可导。对 $F(x)=\frac{f(x)}{x}$、$G(x)=\frac{1}{x}$ 使用柯西中值定理，证明存在 $\xi\in(a,b)$ 使 $\frac{af(b)-bf(a)}{a-b}=f(\xi)-\xi f'(\xi)$。
\end{exercise}
\begin{exercise}{315}{常规·幂函数参照}
设 $0<a<b$，$f\in C[a,b]$，在 $(a,b)$ 可导。对 $f(x)$、$g(x)=x^3$ 使用柯西中值定理，证明存在 $\xi\in(a,b)$ 使 $\frac{f(b)-f(a)}{b-a}=\frac{(a^2+ab+b^2)f'(\xi)}{3\xi^2}$。
\end{exercise}
\begin{exercise}{316}{常规·复合函数中值定理}
对 $F(t)=\sin t$ 在 $[0,1]$ 上使用中值定理，证明存在 $\xi\in(1,e)$ 使 $\sin1=\cos(\ln\xi)$。
\end{exercise}
\begin{exercise}{317}{基础·中值定理变式}
对 $f(x)=e^x$ 在 $[0,1]$ 上使用拉格朗日中值定理，求中值点。
\end{exercise}
\begin{exercise}{318}{常规·罗尔定理综合}
设 $f\in C[0,1]$，在 $(0,1)$ 可导，$f(0)=0$、$f(1)=1$。对 $F(x)=e^x[f(x)-x]$ 使用罗尔定理，证明存在 $\xi\in(0,1)$ 使 $f'(\xi)+f(\xi)-\xi=1$。
\end{exercise}
\subsection{洛必达法则}
\begin{exercise}{319}{基础·零比零型极限}
用洛必达法则求 $\lim_{x\to1}\frac{x^3-1}{x^2-1}$。
\end{exercise}
\begin{exercise}{320}{基础·对数极限}
用洛必达法则求 $\lim_{x\to1}\frac{\ln x}{x-1}$。
\end{exercise}
\begin{exercise}{321}{基础·混合函数极限}
用洛必达法则求 $\lim_{x\to1}\frac{x^3-1+\ln x}{e^x-e}$。
\end{exercise}
\begin{exercise}{322}{基础·无穷比无穷型}
求 $\lim_{x\to+\infty}\frac{x^2}{e^x}$ 与 $\lim_{x\to+\infty}\frac{x^4+e^x}{xe^x}$。
\end{exercise}
\begin{exercise}{323}{常规·连续使用洛必达法则}
求 $\lim_{x\to0}\frac{x-\sin x}{x^2\sin x}$。提示：可先将分母与 $x^3$ 比较，再对 $\frac{x-\sin x}{x^3}$ 使用洛必达法则。
\end{exercise}
\begin{exercise}{324}{基础·对数与幂比较}
求 $\lim_{x\to+\infty}\frac{\ln x}{\sqrt x}$。
\end{exercise}
\begin{exercise}{325}{常规·二阶不定式}
求 $\lim_{x\to0}\frac{e^x-\sin x-1}{1-\sqrt{1-x^2}}$。
\end{exercise}
\begin{exercise}{326}{常规·反正切差极限}
求 $\lim_{x\to0}\frac{\arctan x-x}{x^3}$。
\end{exercise}
\begin{exercise}{327}{基础·乘积极限}
求 $\lim_{x\to0}\frac{\sin x(e^{2x}-1)}{x\ln(1+x)}$。
\end{exercise}
\begin{exercise}{328}{常规·极限反求参数}
求常数 $a$，使 $\lim_{x\to0}\frac{a\tan x+1-\cos x}{\ln(1+2x)}=2$。
\end{exercise}
\begin{exercise}{329}{常规·二阶对数极限}
用洛必达法则求 $\lim_{x\to0}\frac{x-\ln(1+x)}{x^2}$。
\end{exercise}
\begin{exercise}{330}{常规·嵌套三角极限}
求 $\lim_{x\to0}\frac{\sin x-\sin(\sin x)}{x^3}$。提示：令 $u=\sin x$，使用 $u-\sin u\sim\frac{u^3}{6}$。
\end{exercise}
\begin{exercise}{331}{常规·无穷量相减}
求 $\lim_{x\to0}\frac1x(\frac{1}{\sin x}-\frac{1}{\tan x})$。提示：先合并为一个分式。
\end{exercise}
\begin{exercise}{332}{基础·指数衰减}
用洛必达法则求 $\lim_{x\to+\infty}x^3e^{-2x}$。
\end{exercise}
\begin{exercise}{333}{常规·零乘无穷型}
求 $\lim_{x\to1}(1-x)\tan(\frac{\pi x}{2})$。提示：令 $t=1-x$。
\end{exercise}
\begin{exercise}{334}{常规·三角无穷量之差}
求 $\lim_{x\to\frac{\pi}{2}}(\sec x-\tan x)$。提示：乘共轭式或写成 $\frac{\cos x}{1+\sin x}$。
\end{exercise}
\begin{exercise}{335}{常规·倒数差极限}
求 $\lim_{x\to0}[\frac{1}{x}-\frac{1}{e^x-1}]$。提示：先通分。
\end{exercise}
\begin{exercise}{336}{常规·二阶相消}
求 $\lim_{x\to0}[\frac{1}{x^2}-\frac{1}{x\tan x}]$。可使用 $\tan x=x+\frac{x^3}{3}+o(x^3)$。
\end{exercise}
\begin{exercise}{337}{基础·零的零次幂型}
求 $\lim_{x\to0^+}x^x$。提示：先求 $x\ln x$ 的极限。
\end{exercise}
\begin{exercise}{338}{基础·幂指极限}
求 $\lim_{x\to1}x^{\frac{1}{1-x}}$。
\end{exercise}
\begin{exercise}{339}{基础·一的无穷次幂型}
求 $\lim_{x\to0}(1+3x)^{\frac{2}{\sin x}}$，只在底数为正的充分小邻域内考虑。
\end{exercise}
\begin{exercise}{340}{常规·二阶幂指极限}
求 $\lim_{x\to0}(\cos x)^{\frac{1}{\ln(1+x^2)}}$。提示：取对数后使用洛必达法则。
\end{exercise}
\begin{exercise}{341}{基础·无穷的零次幂型}
求 $\lim_{x\to+\infty}x^{\frac{1}{x}}$。
\end{exercise}
\begin{exercise}{342}{常规·指数平均的幂极限}
求 $\lim_{x\to0}[\frac{e^x+e^{3x}}{2}]^{\frac{1}{x}}$。
\end{exercise}
\begin{exercise}{343}{常规·指数差极限}
求 $\lim_{x\to0}\frac{e^{\sin x}-1}{x}$。
\end{exercise}
\begin{exercise}{344}{常规·连续补值参数}
设 $f(x)=\frac{\sin2x+e^{2ax}-1}{x}$（$x\ne0$），$f(0)=a$。求使 $f$ 在 $0$ 处连续的 $a$。
\end{exercise}
\begin{exercise}{345}{基础·幂函数连续补值}
设 $f(x)=(\cos x)^{x^2}$（$0<|x|<\frac{\pi}{2}$），$f(0)=a$。求使 $f$ 在 $0$ 处连续的 $a$。
\end{exercise}
\begin{exercise}{346}{常规·端点相消极限}
求 $\lim_{x\to1^-}[\frac{1}{\sin(\pi x)}-\frac{1}{\pi(1-x)}]$。可令 $t=\pi(1-x)$，使用 $t-\sin t\sim\frac{t^3}{6}$。
\end{exercise}
\begin{exercise}{347}{常规·递推极限}
设 $0<x_1<\frac{\pi}{2}$，$x_{n+1}=\sin x_n$。证明 $x_n$ 单调递减且有下界，求其极限。可使用 $0<\sin t<t$（$0<t<\frac{\pi}{2}$）。
\end{exercise}
\begin{exercise}{348}{基础·无穷远极限变式}
求 $\lim_{x\to+\infty}\frac{x+\sin x}{x}$。提示：使用有界因子估计。
\end{exercise}
\begin{exercise}{349}{常规·对数化极限综合}
求 $\lim_{x\to0}(1+x+2x^2)^{\frac{1}{x}}$，并求其对数的极限。
\end{exercise}
\subsection{泰勒公式}
\begin{exercise}{350}{基础·麦克劳林公式}
写出 $e^x$ 在 $0$ 处的三阶麦克劳林公式，余项用 $o(x^3)$ 表示，并用多项式估算 $e^{0.1}$。
\end{exercise}
\begin{exercise}{351}{常规·泰勒公式求极限}
用泰勒公式求 $\lim_{x\to0}\frac{x\ln(1+x)}{e^x-x-1}$。
\end{exercise}
\begin{exercise}{352}{常规·三阶相消}
用 $\sin x$、$\cos x$ 的三阶以内展开求 $\lim_{x\to0}\frac{\sin x-x\cos x}{x^3}$。
\end{exercise}
\begin{exercise}{353}{常规·对称根式展开}
求 $\lim_{x\to0}\frac{\sqrt{1+x}+\sqrt{1-x}-2}{x^2}$。提示：将两个根式展开到二阶。
\end{exercise}
\begin{exercise}{354}{常规·泰勒系数反求参数}
求常数 $a,b$，使 $\lim_{x\to0}\frac{\ln(1+x)-ax-bx^2}{x^2}=2$。
\end{exercise}
\begin{exercise}{355}{基础·根式二阶展开}
写出 $\sqrt{1+x^2}$ 在 $0$ 处到 $x^2$ 项的展开，并求 $\lim_{x\to0}\frac{\sqrt{1+x^2}-1}{x^2}$。
\end{exercise}
\begin{exercise}{356}{常规·指数与余弦比较}
求 $\lim_{x\to0}\frac{\cos x-e^{\frac{x^2}{2}}}{x^2}$。
\end{exercise}
\begin{exercise}{357}{常规·已知余项求极限}
已知 $\sin6x+xf(x)=o(x^3)$（$x\to0$）。利用 $\sin6x=6x-36x^3+o(x^3)$，求 $\lim_{x\to0}\frac{6+f(x)}{x^2}$。
\end{exercise}
\begin{exercise}{358}{基础·展开式匹配}
求 $a,b$，使 $e^x-(1+bx+ax^2)=o(x^2)$（$x\to0$）。
\end{exercise}
\begin{exercise}{359}{常规·根式差主项}
求 $\sqrt{1-2x}-\sqrt[3]{1-3x}$ 在 $x\to0$ 时的二阶主项。可使用 $(1+u)^\alpha=1+\alpha u+\frac{\alpha(\alpha-1)u^2}{2}+o(u^2)$。
\end{exercise}
\begin{exercise}{360}{基础·对数展开}
设 $f(x)=\ln(1-3x)$。写出 $f$ 在 $0$ 处的三阶麦克劳林公式，并求 $f'''(0)$。
\end{exercise}
\begin{exercise}{361}{常规·拉格朗日余项}
用 $e^x=1+x+\frac{e^\xi x^2}{2}$（$\xi$ 介于 $0$ 与 $x$ 之间）证明 $e^x>1+x$ 对 $x\ne0$ 成立。
\end{exercise}
\begin{exercise}{362}{常规·泰勒公式数值近似}
用 $\sin x$ 的三次麦克劳林多项式估算 $\sin0.1$，并利用四阶拉格朗日余项给出误差上界。
\end{exercise}
\begin{exercise}{363}{常规·二阶泰勒余项}
已知 $f\in C^2[0,1]$，$f(0)=f'(0)=0$、$|f''(x)|\le2$。用泰勒公式证明 $|f(x)|\le x^2$（$0\le x\le1$）。
\end{exercise}
\begin{exercise}{364}{基础·泰勒展开变式}
写出 $\cos2x$ 在 $0$ 处到 $x^2$ 项的展开，余项用 $o(x^2)$ 表示。
\end{exercise}
\begin{exercise}{365}{常规·展开求极限综合}
用泰勒公式求 $\lim_{x\to0}\frac{e^{2x}-1-2x}{1-\cos2x}$。
\end{exercise}
\subsection{函数的单调性与曲线的凹凸性}
\begin{exercise}{366}{基础·单调区间}
求 $f(x)=x-\frac32x^{\frac{2}{3}}$ 的单调区间，其中分数幂按实立方根定义。提示：分别在 $x<0$ 与 $x>0$ 上求导，并检查 $x=0$。
\end{exercise}
\begin{exercise}{367}{常规·导数证明不等式}
用导数证明 $\tan x>x$（$0<x<\frac{\pi}{2}$）。提示：考察 $F(x)=\tan x-x$。
\end{exercise}
\begin{exercise}{368}{常规·对数函数比较}
用导数证明 $\ln(1+x)>\frac{x}{1+x}$（$x>0$）。提示：考察两边之差。
\end{exercise}
\begin{exercise}{369}{基础·单调性与导数}
求 $f(x)=x^3$ 的导数，说明 $f$ 在实轴上的单调性，并计算 $f'(0)$。
\end{exercise}
\begin{exercise}{370}{常规·方程唯一根}
证明 $x+\sqrt x-1=0$ 在 $(0,1)$ 内有唯一实根。提示：使用零点定理和单调性。
\end{exercise}
\begin{exercise}{371}{常规·三次方程根数}
利用 $f(x)=x^3-3x$ 的单调区间，求方程 $x^3-3x=0$ 的实根，并判断方程 $x^3-3x=3$ 的实根个数。
\end{exercise}
\begin{exercise}{372}{常规·高次方程唯一根}
用导数与无穷远极限证明方程 $x^5+x+1=0$ 有唯一实根。
\end{exercise}
\begin{exercise}{373}{常规·商函数单调性}
求 $F(x)=\frac{\ln(1+x)}{x}$（$x>0$）的单调性。可使用 $\ln(1+x)>\frac{x}{1+x}$。
\end{exercise}
\begin{exercise}{374}{常规·指数函数比较}
用导数证明 $(e+x)^e<e^{e+x}$（$x>0$）。提示：取对数后考察 $e+x-e\ln(e+x)$。
\end{exercise}
\begin{exercise}{375}{常规·中值定理与导数下界}
设 $1<a<b<2$。用中值定理证明 $\ln b-\ln a>\frac{b-a}{2}$。
\end{exercise}
\begin{exercise}{376}{基础·凹凸性与拐点}
求 $f(x)=x^3-3x^2$ 的凹凸区间和拐点，注明各区间内 $f''$ 的符号。
\end{exercise}
\begin{exercise}{377}{基础·三角函数凹凸性}
求 $f(x)=1+\sin x$ 在 $(\pi,\frac{3\pi}{2})$ 上的单调性与凹凸性，注明 $f'$、$f''$ 的符号。
\end{exercise}
\begin{exercise}{378}{常规·指数中点不等式}
利用 $(e^{\frac{x}{2}}-e^{\frac{y}{2}})^2\ge0$ 证明 $\frac{e^x+e^y}{2}\ge e^{\frac{x+y}{2}}$，并求等号条件。
\end{exercise}
\begin{exercise}{379}{常规·分段曲线凹凸性}
将 $f(x)=|x^2-1|$ 写为分段函数，求各段的二阶导数和凹凸性，并判断 $x=\pm1$ 处是否可导。
\end{exercise}
\begin{exercise}{380}{基础·拐点条件}
求曲线 $y=2x^3-3x^2+1$ 的拐点坐标，并验证二阶导数在其两侧变号。
\end{exercise}
\begin{exercise}{381}{常规·偶函数导数符号}
设 $f\in C^2(\mathbb R)$ 为偶函数，在 $(-\infty,0)$ 上 $f'>0$、$f''<0$。使用 $f'$ 为奇函数、$f''$ 为偶函数，求它们在 $(0,+\infty)$ 上的符号。
\end{exercise}
\begin{exercise}{382}{基础·水平拐点}
求 $f(x)=x^3$ 的驻点与拐点，并判断驻点是否为极值点。
\end{exercise}
\begin{exercise}{383}{常规·导数关系与凹凸性}
设 $f$ 二阶可导，$f''(x)+[f'(x)]^2=x$ 且 $f'(0)=0$。求 $f''(0)$。若另知 $f'''(0)=1$ 且 $f'''$ 在 $0$ 附近连续，判断 $(0,f(0))$ 是否为拐点。
\end{exercise}
\begin{exercise}{384}{基础·函数与导函数}
设 $f(x)=x^3-3x$。求 $f'$ 的零点及符号，并据此写出 $f$ 的单调区间。
\end{exercise}
\begin{exercise}{385}{基础·单调性变式}
求 $f(x)=x^3-6x$ 的单调区间和驻点。
\end{exercise}
\begin{exercise}{386}{常规·凹凸与零点综合}
设 $f\in C^2[0,1]$，$f(0)=f(1)=0$，且 $f''>0$ 在 $(0,1)$ 成立。先用罗尔定理证明 $f'$ 有零点，再利用 $f'$ 严格递增说明该零点唯一。
\end{exercise}
\subsection{函数的极值和最值}
\begin{exercise}{387}{基础·二阶极值判别}
求 $f(x)=x^3-3x^2+1$ 的驻点，并用二阶导数判别极值。
\end{exercise}
\begin{exercise}{388}{常规·不可导点与极值}
求 $f(x)=2x-\frac32x^{\frac{2}{3}}+1$ 的极值，其中分数幂按实立方根定义。提示：同时检查 $f'=0$ 的点和 $x=0$。
\end{exercise}
\begin{exercise}{389}{基础·指数乘积极值}
求 $f(x)=x2^x$ 的极小值点和极小值。
\end{exercise}
\begin{exercise}{390}{常规·由导函数判别极值}
设 $f'(x)=x(x+2)(x-1)^2(x-3)$。求 $f$ 的极大值点、极小值点，并判断 $x=1$ 是否为极值点。
\end{exercise}
\begin{exercise}{391}{基础·驻点分类}
已知 $f$ 二阶可导，$f'(2)=0$、$f''(2)=3$。判断 $x=2$ 处的极值类型，说明使用的判别法。
\end{exercise}
\begin{exercise}{392}{常规·导数极限与极值}
设 $f$ 二阶可导且 $f'(a)=0$，$\lim_{x\to a}\frac{f'(x)}{x-a}=-1$。求 $f''(a)$，判断 $a$ 处极值类型。
\end{exercise}
\begin{exercise}{393}{常规·增量符号与极值}
设 $\lim_{x\to a}\frac{f(x)-f(a)}{(x-a)^2}=-1$。由极限保号性判断 $a$ 处的极值类型，并求 $f'(a)$。
\end{exercise}
\begin{exercise}{394}{常规·分段函数极值}
设 $f(x)=x^2+1$（$x>0$），$f(x)=xe^x+1$（$x\le0$）。求全部极值点和极值。
\end{exercise}
\begin{exercise}{395}{常规·隐函数极值}
在方程 $x^2+y^2=1$ 的上半圆分支 $y>0$ 上，求 $y$ 关于 $x$ 的极值及取值点。
\end{exercise}
\begin{exercise}{396}{常规·参数与极值}
求 $a$ 的范围，使 $f(x)=x^3+ax$ 在实轴上没有局部极值。提示：考察 $f'(x)=3x^2+a$。
\end{exercise}
\begin{exercise}{397}{基础·闭区间最值}
求 $f(x)=x+\sqrt{1-x}$ 在 $[-5,1]$ 上的最大值和最小值。
\end{exercise}
\begin{exercise}{398}{常规·圆柱体积最值}
将半径为 $1$ 的球内接正圆柱的高记为 $h$（$0<h<2$）。利用 $r^2+(\frac{h}{2})^2=1$ 建立体积函数，求最大体积及对应的高。
\end{exercise}
\begin{exercise}{399}{常规·利润最值}
设总成本 $C(Q)=6Q+200$，总收益 $R(Q)=12Q-0.01Q^2$，$Q\ge0$ 为连续产量。求利润最大时的产量和最大利润。
\end{exercise}
\begin{exercise}{400}{常规·收入成本与最值}
设 $R(x)=26x-2x^2$、$C(x)=8x+x^2$，$x\ge0$。求利润函数及最大利润对应的产量。
\end{exercise}
\begin{exercise}{401}{常规·最值证明不等式}
利用 $F(x)=\frac{x^2}{2}+\frac{1}{2}-x$ 的最小值证明 $\frac{x^2}{2}+\frac{1}{2}\ge x$，并求等号条件。
\end{exercise}
\begin{exercise}{402}{常规·凸性与最小值}
设 $f\in C^2(\mathbb R)$，$f(0)=0$、$f'(0)=1$ 且 $f''>0$。考察 $F(x)=f(x)-x$ 的导数，证明 $f(x)\ge x$，求等号条件。
\end{exercise}
\begin{exercise}{403}{基础·闭区间最值变式}
求 $f(x)=x^2-x$ 在 $[0,1]$ 上的最大值与最小值及取值点。
\end{exercise}
\begin{exercise}{404}{常规·无盖箱体最值}
从边长为 $6$ 的正方形薄板四角各剪去边长为 $x$ 的小正方形，折成无盖长方体。写出体积函数与可行域，求最大体积及对应的 $x$。
\end{exercise}
\subsection{函数图形的描绘}
\begin{exercise}{405}{基础·水平与铅直渐近线}
求 $y=\frac{x+1}{x}$ 的水平渐近线和铅直渐近线。
\end{exercise}
\begin{exercise}{406}{基础·有理函数渐近线}
求 $y=\frac{1}{1-x^2}$ 的水平渐近线和铅直渐近线。
\end{exercise}
\begin{exercise}{407}{常规·可去间断与渐近线}
求 $f(x)=\frac{x^2-x-2}{x^2-1}$ 的水平、铅直渐近线，并说明 $x=-1$ 是否为铅直渐近线。
\end{exercise}
\begin{exercise}{408}{基础·斜渐近线}
用多项式除法求 $y=\frac{x^2}{x-1}$ 的斜渐近线和铅直渐近线。
\end{exercise}
\begin{exercise}{409}{常规·指数复合渐近线}
求 $y=(2x-1)e^{\frac{1}{x}}$ 的斜渐近线及 $x\to0^+$ 的极限。可使用 $e^{\frac{1}{x}}=1+\frac{1}{x}+o(\frac{1}{x})$（$|x|\to\infty$）。
\end{exercise}
\begin{exercise}{410}{基础·振荡因子与渐近线}
设 $f(x)=x\sin(\frac{1}{x})$（$x>0$）。求 $x\to+\infty$ 时的水平渐近线，并求 $x\to0^+$ 的极限。
\end{exercise}
\begin{exercise}{411}{常规·指数曲线渐近线}
求 $f(x)=xe^{\frac{1}{x}}$ 的斜渐近线及铅直渐近线。提示：分别计算 $x\to0^\pm$ 的极限。
\end{exercise}
\begin{exercise}{412}{常规·对数函数水平渐近线}
分别求 $f(x)=\frac{\ln(1+e^x)}{x}$ 在正、负无穷方向的水平渐近线。
\end{exercise}
\begin{exercise}{413}{基础·反三角函数渐近线}
求 $f(x)=\arctan x$ 在正、负无穷方向的水平渐近线，并求它与坐标轴的交点。
\end{exercise}
\begin{exercise}{414}{常规·有理函数图形分析}
对 $f(x)=\frac{x^2}{x-1}$，求定义域、渐近线、单调区间和极值，据此画出图形草图。
\end{exercise}
\begin{exercise}{415}{常规·周期函数图形分析}
求 $f(x)=\cos x$ 的最小正周期，在 $[-\pi,\pi]$ 上求单调区间和极值，并画出该区间的图形草图。
\end{exercise}
\begin{exercise}{416}{基础·斜渐近线变式}
求 $f(x)=x+\frac{\sin x}{x}$（$x\ne0$）的斜渐近线，并求 $x\to0$ 时的极限。
\end{exercise}
\begin{exercise}{417}{常规·渐近线与单调性综合}
求 $f(x)=\frac{x^2+1}{x-1}$ 的铅直、斜渐近线及单调区间。
\end{exercise}
\subsection{曲率}
\begin{exercise}{418}{基础·显函数曲率}
用 $K=\frac{|y''|}{(1+(y')^2)^{\frac{3}{2}}}$ 求曲线 $y=\frac{1}{x}$ 在 $(1,1)$ 处的曲率与曲率半径。
\end{exercise}
\begin{exercise}{419}{常规·参数曲线曲率}
求曲线 $x=2\cos^3t$、$y=2\sin^3t$ 在 $t=\frac{\pi}{4}$ 处的曲率半径。可使用 $K=\frac{|x'y''-y'x''|}{((x')^2+(y')^2)^{\frac{3}{2}}}$。
\end{exercise}
\begin{exercise}{420}{常规·曲率与曲率圆}
求抛物线 $y=x^2$ 在原点处的曲率和曲率半径，并写出曲率圆方程。提示：圆心在 $y$ 轴正半轴上。
\end{exercise}
\begin{exercise}{421}{基础·根式曲线曲率}
求 $y=\sqrt x$ 在 $(1,1)$ 处的曲率半径。
\end{exercise}
\begin{exercise}{422}{基础·抛物线曲率变式}
求 $y=2x^2$ 在原点处的曲率及曲率半径。
\end{exercise}
\begin{exercise}{423}{常规·椭圆曲率综合}
设椭圆 $x=2\cos t$、$y=\sin t$。求 $t=0$ 与 $t=\frac{\pi}{2}$ 两点的曲率半径并比较大小。可使用参数曲线曲率公式。
\end{exercise}
